\subsection*{(f) }
\addcontentsline{toc}{subsection}{(f) }
Pelo item anterior, temos que
\begin{equation*}
    L_t(\hat{\alpha}_t, \hat{\theta}_t) \le \sqrt{( 1 - 4 \gamma^2 )} L_{t-1}(\hat{\alpha}_{t-1}, \hat{\theta}_{t-1})
\end{equation*}
Extendendo para t = 0, temos que
\begin{equation*}
    L_t(\hat{\alpha}_t, \hat{\theta}_t) \le ( 1 - 4 \gamma^2 )^{t/2} L_0(\hat{\alpha}_0, \hat{\theta}_0)
\end{equation*}
Como $L_0(\hat{\alpha}_0, \hat{\theta}_0) = 1$, temos que
\begin{equation*}
    L_t(\hat{\alpha}_t, \hat{\theta}_t) \le ( 1 - 4 \gamma^2 )^{t/2}
\end{equation*}

Como $1+x \le e^x$, temos que
\begin{equation*}
    1-4\gamma^2 \le e^{-4\gamma^2}
\end{equation*}

\begin{equation*}
    (1-4\gamma^2)^{t/2} \le e^{-2\gamma^2 t}
\end{equation*}
\begin{equation*}
    L_t(\hat{\alpha}_t, \hat{\theta}_t) \le e^{-2 \gamma^2 t}
\end{equation*}
